% =============================================================================
% 測圓海鏡 (Sea Mirror of Circle Measurements) — Volumes 1–3
% Author: Li Ye (Li Zhi, 1192–1279)
% Typeset in LaTeX with XeLaTeX + xeCJK
% =============================================================================
\documentclass[11pt,a4paper]{article}

% -------- Core Packages --------
\usepackage{fontspec}
\usepackage{polyglossia}
\setmainlanguage{english}
\usepackage{amsmath,amssymb,amsthm}
\usepackage{geometry}
\usepackage{graphicx}
\usepackage{xcolor}
\usepackage{titlesec}
\usepackage{fancyhdr}
\usepackage{enumitem}
\usepackage{array,booktabs,longtable,multirow}
\usepackage{hyperref}
\usepackage{tocloft}
\usepackage{indentfirst}
\usepackage{xeCJK}

% -------- CJK Font Setup --------
\setCJKmainfont{NotoSerifCJKsc-Regular.otf}[
  Path=fonts/noto-sc-extracted/,
  BoldFont=NotoSerifCJKsc-Bold.otf
]
\setCJKsansfont{NotoSansCJKsc-Regular.otf}[
  Path=fonts/noto-sc-extracted/,
  BoldFont=NotoSansCJKsc-Bold.otf
]
\setCJKmonofont{NotoSansCJKsc-Regular.otf}[
  Path=fonts/noto-sc-extracted/,
  BoldFont=NotoSansCJKsc-Bold.otf
]

% -------- Latin Font Setup --------
\setmainfont{Linux Libertine O}
\setsansfont{Linux Biolinum O}

% -------- Page Geometry --------
\geometry{
  paperwidth=210mm,paperheight=297mm,
  left=25mm,right=25mm,top=30mm,bottom=30mm,
  headheight=14pt,footskip=30pt
}

% -------- Colors --------
\definecolor{titlecolor}{RGB}{45,55,72}
\definecolor{sectioncolor}{RGB}{55,65,81}
\definecolor{notecolor}{RGB}{120,53,15}
\definecolor{rulecolor}{RGB}{180,160,140}
\definecolor{prefacecolor}{RGB}{80,60,40}
\definecolor{problemcolor}{RGB}{30,60,90}

% -------- Section Formatting --------
\titleformat{\section}
  {\normalfont\Large\bfseries\color{sectioncolor}}
  {\thesection}{1em}{}
\titleformat{\subsection}
  {\normalfont\large\bfseries\color{sectioncolor}}
  {\thesubsection}{1em}{}
\titleformat{\subsubsection}
  {\normalfont\normalsize\bfseries\color{sectioncolor}}
  {\thesubsubsection}{1em}{}

% -------- Header/Footer --------
\pagestyle{fancy}
\fancyhf{}
\fancyhead[L]{\small\textit{測圓海鏡 — Sea Mirror of Circle Measurements}}
\fancyhead[R]{\small\thepage}
\renewcommand{\headrulewidth}{0.4pt}
\renewcommand{\footrulewidth}{0pt}

% -------- Custom Commands --------
\newcommand{\zhsection}[1]{\section*{#1}\addcontentsline{toc}{section}{#1}}
\newcommand{\zhsubsection}[1]{\subsection*{#1}\addcontentsline{toc}{subsection}{#1}}
\newcommand{\zhsubsubsection}[1]{\subsubsection*{#1}\addcontentsline{toc}{subsubsection}{#1}}

% Problem/Solution structure
\newcommand{\problem}[1]{\par\noindent\textbf{\color{problemcolor}#1}\par\nopagebreak}
\newcommand{\answer}[1]{\par\noindent\textbf{答} #1\par}
\newcommand{\method}[1]{\par\noindent\textbf{法曰} #1\par}
\newcommand{\workings}[1]{\par\noindent\textbf{草曰} #1\par}
\newcommand{\remark}[1]{\par\noindent\textit{〔按〕}#1\par}

% Chinese terms
\newcommand{\term}[1]{\textbf{#1}}

% Hyperref
\hypersetup{
  colorlinks=true,
  linkcolor=sectioncolor,
  citecolor=sectioncolor,
  urlcolor=sectioncolor,
  pdfencoding=auto
}

\setcounter{tocdepth}{2}

% =============================================================================
\begin{document}

% -------- Title Page --------
\begin{titlepage}
\centering
\vspace*{2cm}

{\fontsize{36}{44}\selectfont\CJKfamily{sf}\color{titlecolor}
測圓海鏡}

\vspace{0.8cm}

{\fontsize{18}{22}\selectfont\color{sectioncolor}
Sea Mirror of Circle Measurements}

\vspace{0.5cm}

{\large\color{rulecolor}\rule{0.6\textwidth}{0.4pt}}

\vspace{0.8cm}

{\LARGE 卷一 — 卷三}

{\large Volumes I – III}

\vspace{1.5cm}

{\large\color{sectioncolor}
元 \quad 李冶 \quad 撰}

\vspace{0.3cm}

{\normalsize\color{sectioncolor}
(Li Ye / Li Zhi, 1192–1279)}

\vspace{2cm}

{\color{rulecolor}\rule{0.4\textwidth}{0.4pt}}

\vspace{1cm}

{\normalsize
\textit{以天元術解勾股容圓一百七十問}}

\vspace{0.3cm}

{\small\color{notecolor}
\textit{Using the Method of the Celestial Element (Tian Yuan Shu)}\\
\textit{to Solve 170 Problems on the Inscribed Circle}}

\vspace{1.5cm}

{\small
A foundational text of medieval Chinese algebra,\\
completed in the year 1248 of the Common Era.
}

\vfill

{\small Typeset in \LaTeX\ with XeLaTeX + xeCJK}

\end{titlepage}

% -------- Table of Contents --------
\tableofcontents
\newpage

% =============================================================================
\section{Introduction}
\label{sec:intro}
\addcontentsline{toc}{section}{Introduction}

The \textit{Ceyuan Haijing} (測圓海鏡, literally ``Sea Mirror of Circle Measurements'') is the masterpiece of Li Ye (李冶, also known as Li Zhi, 1192–1279), one of the four great mathematicians of the Song--Yuan period of China, alongside Yang Hui, Qin Jiushao, and Zhu Shijie.

Completed in 1248, this work systematically applies the \textbf{tian yuan shu} (天元術, ``method of the celestial element'')---an algebraic method for setting up and solving polynomial equations---to a total of 170 geometric problems concerning a circle inscribed in or related to a right triangle (勾股容圓).

The book is organized as follows:
\begin{itemize}[leftmargin=2em]
  \item \textbf{Volume 1} establishes the theoretical foundation: the \textit{Circle City Diagram} (圓城圖式), the \textit{General Rate Names} (總率名號), the \textit{Current Question Numbers} (今問正數), and the \textit{Identification Miscellaneous Records} (識別雜記) containing 692 geometric formulas.
  \item \textbf{Volumes 2--12} present 170 problems, solved by the tian yuan shu method. The problems involve equations of degrees ranging from 1 to 6.
\end{itemize}

The present document contains the complete text of \textbf{Volumes 1–3}, comprising all foundational material and the first suite of problems.

\vspace{0.5cm}
\noindent\textcolor{rulecolor}{\rule{\textwidth}{0.4pt}}

% =============================================================================
\newpage
\section{Original Preface (原序)}
\label{sec:preface}

{\color{prefacecolor}
\begin{quote}
數本難窮，吾欲以力強窮之，彼其數不惟不能得其凡，而吾之力且憊矣。然則數果不可以窮耶？既已名之數矣，則又何為而不可窮也。

故謂數為難窮斯可，謂數為不可窮斯不可。何則？彼其冥冥之中，固有昭昭者存，夫昭昭者其自然之數也。非自然之數，其自然之理也。數一出於自然，吾欲以力強窮之，使隸首復生亦末如之何也已。

苟能推自然之理，以明自然之數，則雖遠而乾端坤倪，幽而神情鬼狀，未有不合者矣。
\end{quote}
}

\vspace{0.5cm}
\noindent\textit{Translation:} Numbers are fundamentally difficult to exhaust; if I try to force them by effort, not only will I fail to apprehend their essence, but my strength will be depleted as well. Yet can numbers truly not be exhausted? Since they are already called \textit{numbers}, why then could they not be exhausted?

It is acceptable to say that numbers are difficult to exhaust; but it is unacceptable to say they cannot be exhausted. Why? In the midst of darkness, there is indeed brightness. That brightness is the natural number. It is not merely the natural number---it is the natural principle. Since numbers arise from nature, if I try to force them by effort, even if Li Shou (the legendary inventor of arithmetic) were reborn, he could do no better.

If one can extend natural principles to illuminate natural numbers, then even things as distant as the corners of heaven and earth, or as mysterious as spirits and ghosts, will all be in agreement.

\vspace{0.5cm}

{\color{prefacecolor}
\begin{quote}
予自幼喜算數，恆病夫考圓之術例出於牽強，殊乖於自然，如古率、徽率、密率之不同，截弧、截矢、截背之互見，內外諸角，析剖支條，莫不各自名家，與世作法。及反覆研究，卒無以當吾心焉。

老大以來，得洞淵九容之說，日夕玩繹，而向之病我者，使爆然落去而無遺余。山中多暇，客有從余求其說者，於是乎又為衍之，遂累一百七十問。

既成編，客復目之《測圓海鏡》，蓋取夫天臨海鏡之義也。
\end{quote}
}

\vspace{0.5cm}
\noindent\textit{Translation:} Since youth I have delighted in arithmetic. I was always troubled that the methods for investigating circles seemed forced and artificial, contrary to nature---the ancient rate, Hui's rate, and the close rate all differing; the arc, sagitta, and back each being considered separately; the internal and external angles analyzed into minute subdivisions---each school establishing its own name and method for the world. After repeated study, none could satisfy my mind.

In old age, I obtained the \textit{Dongyuan Nine Containments} (洞淵九容) theory, and pondered it day and night. What had previously troubled me suddenly fell away completely, leaving no trace. With leisure time in the mountains, when guests asked me to explain this theory, I expanded upon it, eventually accumulating \textbf{one hundred and seventy problems}.

When the compilation was complete, a guest named it \textit{Ceyuan Haijing}---taking the meaning of ``heaven overlooking the sea mirror.''

\vspace{0.5cm}

{\color{prefacecolor}
\begin{quote}
昔半山老人集《唐百家詩選》，自謂廢日力於此，良可惜。明道先生以上蔡謝君記誦為玩物喪志。夫文史尚矣，猶之為不足貴，況九九賤技能乎！

嗜好酸鹹，平生每痛自戒勅，竟莫能已，類有物慿之者，吾亦不知其然而然也。

故嘗私為之解曰：由技兼於事者言之，夷之禮，夔之樂，亦不免為一技。由技進乎道者言之，石之斤，扁之輪，非聖人之所與乎。

覽吾之編，察吾苦心，其憫我者當百數，其笑我者當千數。乃若吾之所自得，則自得焉耳，寧復為人憫笑計哉。

\hfill 李冶序
\end{quote}
}

\newpage

% =============================================================================
\section{Volume 1 \quad 卷一}
\label{sec:vol1}

Volume 1 is the theoretical foundation of the entire work. It comprises four major sections:

\begin{enumerate}[leftmargin=2em]
  \item \textbf{圓城圖式} (Circle City Diagram) --- the master geometric configuration
  \item \textbf{總率名號} (General Rate Names) --- definitions of all 15 right triangles and their elements
  \item \textbf{今問正數} (Current Question Numbers) --- numerical values for all line segments
  \item \textbf{識別雜記} (Identification Miscellaneous Records) --- 692 geometric formulas
\end{enumerate}

% -----------------------------------------------------------------------------
\subsection{Circle City Diagram (圓城圖式)}
\label{subsec:diagram}

The \textit{Circle City Diagram} is the master geometric configuration upon which all 170 problems are based. It consists of a large right triangle (called 勾股形, ``leg--leg shape'') with its inscribed circle, together with specific points and lines.

The scenario Li Ye constructs is:

\begin{quote}
假令有圓城一所，不知周徑，四面開門，門外縱橫各有十字大道。其西北十字道頭定為乾地，其東北十字道頭定為艮地，其東南十字道頭定為巽地，其西南十字道頭定為坤地。所有測望雜法一一設問如後。
\end{quote}

\textit{Translation:} Suppose there is a circular city whose circumference and diameter are unknown. It has gates on all four sides, and outside each gate there are crossroads running north--south and east--west. The northwest crossroad is designated \textbf{Qian} (乾), the northeast \textbf{Gen} (艮), the southeast \textbf{Xun} (巽), and the southwest \textbf{Kun} (坤). All surveying methods are set forth as questions below.

\vspace{0.3cm}

\textbf{The 22 Named Points.} The largest right triangle has vertices \textbf{Tian} (天, Heaven), \textbf{Di} (地, Earth), and \textbf{Qian} (乾). The center of the inscribed circle is called \textbf{Xin} (心, Heart). Key points include intersections of horizontal and vertical lines through the center and other constructed lines. Each point is designated by a single Chinese character. The complete list of named points is as follows:

\begin{center}
\begin{tabular}{llllllll}
\toprule
天 & 地 & 乾 & 心 & 日 & 南 & 北 & 川 \\
東 & 西 & 艮 & 坤 & 山 & 月 & 巽 & 青 \\
泉 & 朱 & 金 & 泛 & 旦 & 夕 & & \\
\bottomrule
\end{tabular}
\end{center}

Each of these twenty-two characters denotes a specific point in the geometric configuration. The lines connecting them form fifteen distinct right triangles, each sharing a common relationship with the inscribed circle of radius 120 paces.

\vspace{0.3cm}

\textbf{The Geometric Configuration.} The diagram is constructed as follows. The great right triangle 天地乾 has its right angle at 乾. The circle is inscribed within this triangle, tangent to all three sides. The center of the circle is 心. Vertical and horizontal lines through 心 intersect the sides of the triangle and each other at key points, generating the remaining nineteen named points.

The city gates are located at the four cardinal points where the circle meets the horizontal and vertical axes through 心: 北 (North Gate), 南 (South Gate), 東 (East Gate), and 西 (West Gate). The crossroads outside each gate are named according to the Eight Trigrams: 乾 (northwest), 艮 (northeast), 巽 (southeast), and 坤 (southwest).

\vspace{0.3cm}

\textbf{Fundamental Relationships.} The diagram embodies the classical relationship between a right triangle and its inscribed circle. For the great triangle 天地乾:
\[
\text{勾} + \text{股} = \text{弦} + \text{圓徑}
\]
That is, the sum of the two legs equals the sum of the hypotenuse and the diameter of the inscribed circle. This is one of the 692 formulas catalogued in the \textit{Identification Miscellaneous Records}.

% -----------------------------------------------------------------------------
\subsection{General Rate Names (總率名號)}
\label{subsec:names}

The work studies \textbf{15 right triangles}, all having a segment of the \textit{Tiandi} (天地) line as their hypotenuse (弦, ``string''). The \textit{General Rate Names} define these triangles and their three sides.

In each triangle:
\begin{itemize}[leftmargin=2em]
  \item \textbf{弦} ($c$) --- the hypotenuse (斜邊)
  \item \textbf{股} ($b$) --- the longer vertical leg (縱邊)
  \item \textbf{勾} ($a$) --- the shorter horizontal leg (橫邊)
\end{itemize}

\vspace{0.3cm}

The fifteen triangles are generated by the intersection of various lines in the Circle City Diagram. Each triangle is named according to its geometric position or character:

\begin{longtable}{|c|l|l|c|c|c|}
\hline
\textbf{No.} & \textbf{Name} & \textbf{Vertices} & \textbf{弦} & \textbf{股} & \textbf{勾} \\
\hline
\endfirsthead
\hline
\textbf{No.} & \textbf{Name} & \textbf{Vertices} & \textbf{弦} & \textbf{股} & \textbf{勾} \\
\hline
\endhead
1 & 通 (Tong/General) & 天--地--乾 & 通弦 & 通股 & 通勾 \\
2 & 邊 (Bian/Side) & 天--西--川 & 邊弦 & 邊股 & 邊勾 \\
3 & 底 (Di/Base) & 日--地--北 & 底弦 & 底股 & 底勾 \\
4 & 黃廣 (Huangguang) & 天--山--金 & 黃廣弦 & 黃廣股 & 黃廣勾 \\
5 & 黃長 (Huangchang) & 月--地--泉 & 黃長弦 & 黃長股 & 黃長勾 \\
6 & 上高 (Shanggao) & 天--日--旦 & 上高弦 & 上高股 & 上高勾 \\
7 & 下高 (Xiagao) & 日--山--朱 & 下高弦 & 下高股 & 下高勾 \\
8 & 上平 (Shangping) & 月--川--青 & 上平弦 & 上平股 & 上平勾 \\
9 & 下平 (Xiaping) & 川--地--夕 & 下平弦 & 下平股 & 下平勾 \\
10 & 大差 (Dacha) & 天--月--坤 & 大差弦 & 大差股 & 大差勾 \\
11 & 小差 (Xiaocha) & 山--地--艮 & 小差弦 & 小差股 & 小差勾 \\
12 & 皇極 (Huangji) & 日--川--心 & 皇極弦 & 皇極股 & 皇極勾 \\
13 & 太虛 (Taixu) & 月--山--泛 & 太虛弦 & 太虛股 & 太虛勾 \\
14 & 明 (Ming) & 日--月--南 & 明弦 & 明股 & 明勾 \\
15 & 叀 (Zhuan) & 山--川--東 & 叀弦 & 叀股 & 叀勾 \\
\hline
\end{longtable}

Note: 上高 = 下高 and 上平 = 下平, so among 15 triangles, only 13 are distinct.

\vspace{0.3cm}

\textbf{Explanation of Naming.}

\begin{itemize}[leftmargin=2em]
  \item \textbf{通} (General): The largest triangle, containing all others.
  \item \textbf{邊} (Side): Adjacent to the side of the great triangle.
  \item \textbf{底} (Base): Adjacent to the base of the great triangle.
  \item \textbf{黃廣} (Yellow--Wide): Named for its relation to the diameter as the ``yellow'' (central) and ``wide'' extent.
  \item \textbf{黃長} (Yellow--Long): The complement to the Yellow--Wide triangle.
  \item \textbf{上高} and \textbf{下高} (Upper/Lower Height): Derived from the upper and lower portions of the height line.
  \item \textbf{上平} and \textbf{下平} (Upper/Lower Level): Derived from the upper and lower portions of the horizontal line.
  \item \textbf{大差} (Large Difference) and \textbf{小差} (Small Difference): Related to the difference between hypotenuse and legs.
  \item \textbf{皇極} (Imperial Ultimate): The central triangle passing through the circle's center.
  \item \textbf{太虛} (Great Void): The smallest triangle, furthest from the center.
  \item \textbf{明} (Bright) and \textbf{叀} (Special): Two additional triangles generated by interior lines.
\end{itemize}

% -----------------------------------------------------------------------------
\subsection{Current Question Numbers (今問正數)}
\label{subsec:numbers}

This section gives the numerical value of every line segment in the diagram. The \textbf{radius of the inscribed circle is taken as 120 steps} (步), which serves as the standard unit. From this, all other quantities are derived.

The entries are organized by triangle. For each triangle, thirteen numerical quantities are given:

\begin{center}
\begin{tabular}{llll}
弦 & 勾 & 股 & \\
勾股和 & 勾股較 & 勾弦和 & 勾弦較 \\
股弦和 & 股弦較 & 弦較和 & 弦較較 \\
弦和和 & 弦和較 & & \\
\end{tabular}
\end{center}

\vspace{0.3cm}

\textbf{1. 通 (Tong / General)}

\begin{center}
\begin{tabular}{ll}
弦六百八十 & 勾三百二十 \quad 股六百 \\
勾股和九百二十 & 較二百八十 \\
勾和一千 & 較三百六十 \\
股和一千二百八十 & 較八十 \\
較和九百六十 & 較四百 \\
和和一千六百 & 較二百四十 \\
\end{tabular}
\end{center}

\textit{That is:} hypotenuse $c = 680$, short leg $a = 320$, long leg $b = 600$.

\vspace{0.3cm}

\textbf{2. 邊 (Bian / Side)}

\begin{center}
\begin{tabular}{ll}
弦五百四十四 & 勾二百五十六 \quad 股四百八十 \\
勾股和七百三十六 & 較二百二十四 \\
勾和八百 & 較二百八十八 \\
股和一千零二十四 & 較六十四 \\
較和七百六十八 & 較三百二十 \\
和和一千二百八十 & 較一百九十二 \\
\end{tabular}
\end{center}

\vspace{0.3cm}

\textbf{3. 底 (Di / Base)}

\begin{center}
\begin{tabular}{ll}
弦四百二十五 & 勾二百 \quad 股三百七十五 \\
勾股和五百七十五 & 較一百七十五 \\
勾和六百二十五 & 較二百二十五 \\
股和八百 & 較五十 \\
較和六百 & 較二百五十 \\
和和一千 & 較一百五十 \\
\end{tabular}
\end{center}

\vspace{0.3cm}

\textbf{4. 黃廣 (Huang Guang)}

\begin{center}
\begin{tabular}{ll}
弦五百一十 & 勾二百四十 \quad 股四百五十 \\
\multicolumn{2}{c}{(即城徑也 --- this is the city diameter)} \\
勾股和六百九十 & 較二百一十 \\
\end{tabular}
\end{center}

\vspace{0.3cm}

\textbf{5. 黃長 (Huang Chang)}

\begin{center}
\begin{tabular}{ll}
弦二百七十二 & 勾一百二十八 \quad 股二百四十 \\
勾股和三百六十八 & 較一百一十二 \\
勾和四百 & 較一百四十四 \\
股和五百一十二 & 較三十二 \\
較和三百八十四 & 較一百六十 \\
和和六百四十 & 較九十六 \\
\end{tabular}
\end{center}

\vspace{0.3cm}

\textbf{6. 上高 (Shang Gao)}

\begin{center}
\begin{tabular}{ll}
弦二百五十五 & 勾一百二十 \quad 股二百二十五 \\
勾股和三百四十五 & 較一百零五 \\
勾和三百七十五 & 較一百三十五 \\
股和四百八十 & 較三十 \\
較和三百六十 & 較一百五十 \\
和和六百 & 較九十 \\
\end{tabular}
\end{center}

\vspace{0.3cm}

\textbf{7. 下高 (Xia Gao)}

同於上高。

\vspace{0.3cm}

\textbf{8. 上平 (Shang Ping)}

\begin{center}
\begin{tabular}{ll}
弦一百七十 & 勾八十 \quad 股一百五十 \\
勾股和二百三十 & 較七十 \\
勾和二百五十 & 較九十 \\
股和三百二十 & 較二十 \\
較和二百四十 & 較一百 \\
和和四百 & 較六十 \\
\end{tabular}
\end{center}

\vspace{0.3cm}

\textbf{9. 下平 (Xia Ping)}

同於上平。

\vspace{0.3cm}

\textbf{10. 大差 (Da Cha)}

\begin{center}
\begin{tabular}{ll}
弦三百四十 & 勾一百六十 \quad 股三百 \\
勾股和四百六十 & 較一百四十 \\
勾和五百 & 較一百八十 \\
股和六百四十 & 較四十 \\
較和四百八十 & 較二百 \\
和和八百 & 較一百二十 \\
\end{tabular}
\end{center}

\vspace{0.3cm}

\textbf{11. 小差 (Xiao Cha)}

\begin{center}
\begin{tabular}{ll}
弦一百三十六 & 勾六十四 \quad 股一百二十 \\
勾股和一百八十四 & 較五十六 \\
勾和二百 & 較七十二 \\
股和二百五十六 & 較一十六 \\
較和一百九十二 & 較八十 \\
和和三百二十 & 較四十八 \\
\end{tabular}
\end{center}

\vspace{0.3cm}

\textbf{12. 皇極 (Huang Ji)}

\begin{center}
\begin{tabular}{ll}
弦二百五十五 & 勾一百二十 \quad 股二百二十五 \\
勾股和三百四十五 & 較一百零五 \\
勾和三百七十五 & 較一百三十五 \\
股和四百八十 & 較三十 \\
較和三百六十 & 較一百五十 \\
和和六百 & 較九十 \\
\end{tabular}
\end{center}

\vspace{0.3cm}

\textbf{13. 太虛 (Tai Xu)}

\begin{center}
\begin{tabular}{ll}
弦一百零二 & 勾四十八 \quad 股九十 \\
勾股和一百三十八 & 較四十二 \\
勾和一百五十 & 較五十四 \\
股和一百九十二 & 較一十二 \\
較和一百四十四 & 較六十 \\
和和二百四十 & 較三十六 \\
\end{tabular}
\end{center}

\vspace{0.3cm}

\textbf{14. 明 (Ming)}

\begin{center}
\begin{tabular}{ll}
弦一百七十 & 勾八十 \quad 股一百五十 \\
勾股和二百三十 & 較七十 \\
勾和二百五十 & 較九十 \\
股和三百二十 & 較二十 \\
較和二百四十 & 較一百 \\
和和四百 & 較六十 \\
\end{tabular}
\end{center}

\vspace{0.3cm}

\textbf{15. 叀 (Zhuan)}

\begin{center}
\begin{tabular}{ll}
弦六十八 & 勾三十二 \quad 股六十 \\
勾股和九十二 & 較二十八 \\
勾和一百 & 較三十六 \\
股和一百二十八 & 較八 \\
較和九十六 & 較四十 \\
和和一百六十 & 較二十四 \\
\end{tabular}
\end{center}

\vspace{0.3cm}

Each of the 15 triangles yields 13 numerical quantities (弦, 勾, 股, and their ten combinations), giving a total of $15 \times 13 = 195$ numerical entries in the \textit{Current Question Numbers}.

% -----------------------------------------------------------------------------
\subsection{Identification Miscellaneous Records (識別雜記)}
\label{subsec:identification}

This is the theoretical heart of the work. It contains \textbf{692 geometric formulas} relating the various line segments in the diagram. These formulas are purely geometric theorems, independent of any particular numerical values. The problems in all subsequent volumes are solved by using these formulas together with the tian yuan shu.

The section is divided into eight parts:

\begin{enumerate}[leftmargin=2em]
  \item \textbf{諸雜名目} (Various Names) --- general definitions and 30+ theorems
  \item \textbf{五和五較} (Five Sums and Five Differences)
  \item \textbf{諸弦} (Various Hypotenuses)
  \item \textbf{大小差} (Large and Small Differences)
  \item \textbf{諸差} (Various Differences)
  \item \textbf{諸率互見} (Mutual Relations of Rates)
  \item \textbf{四位拾遺} (Four-Place Supplement)
  \item \textbf{拾遺} (Supplement)
\end{enumerate}

\vspace{0.3cm}

\textbf{Sample Definitions from 諸雜名目:}

\begin{center}
\begin{tabular}{ll}
\toprule
\textbf{Name} & \textbf{Definition} \\
\midrule
內率 & 明勾 $\times$ 叀股 \\
外率 & 明股 $\times$ 叀勾 \\
虛率 & 虛勾 $\times$ 虛股 \\
角差 & 高股 $-$ 平勾 \\
次差 & (明股$-$明勾) $+$ (叀股$-$叀勾) \\
混同和 & 黃長股 $+$ 黃廣勾 \\
傍差 & (明股$-$明勾) $-$ (叀股$-$叀勾) \\
\bottomrule
\end{tabular}
\end{center}

\vspace{0.3cm}

\textbf{Fundamental Theorems from 諸雜名目:}

\begin{enumerate}[leftmargin=2em]
  \item 凡大差小差相乘為半段徑冪。\\
  \textit{The product of the large difference and small difference equals half the square of the diameter.}
  
  \item 大差勾小差股相乘同上。\\
  \textit{The product of the large-difference short-leg and the small-difference long-leg equals the same.}
  
  \item 黃廣股黃長勾相乘為經冪。\\
  \textit{The product of the Huangguang long-leg and Huangchang short-leg equals the square of the diameter.}
  
  \item 勾股和即弦黃和。\\
  \textit{The sum of the two legs equals the sum of the hypotenuse and the diameter of the inscribed circle.}\\
  That is: $a + b = c + d$, where $d$ is the diameter of the inscribed circle.
  
  \item 虛勾虛股相乘得半徑冪。\\
  \textit{The product of the Taixu short-leg and long-leg equals the square of the radius.}
  
  \item 明勾明股相乘得半虛徑冪。\\
  \textit{The product of the Ming short-leg and long-leg equals half the square of the Taixu diameter.}
  
  \item 叀勾叀股相乘亦得半虛徑冪。\\
  \textit{The product of the Zhuan short-leg and long-leg also equals half the square of the Taixu diameter.}
  
  \item 皇極勾乘皇極股為圓徑冪。\\
  \textit{The product of the Huangji short-leg and long-leg equals the square of the circle's diameter.}
  
  \item 明弦叀弦相乘為半徑冪。\\
  \textit{The product of the Ming hypotenuse and Zhuan hypotenuse equals the square of the radius.}
  
  \item 虛弦謂之虛黃，即大差差與小差差之共數也。\\
  \textit{The Taixu hypotenuse is called the ``Taixu yellow''; it equals the sum of the large-difference difference and the small-difference difference.}
\end{enumerate}

\vspace{0.3cm}

\textbf{Sample Formulas from 五和五較:}

\begin{itemize}[leftmargin=2em]
  \item 通勾上容圓之徑，即虛勾虛股之和也。
  \item 通股上容圓之徑，即虛勾虛股之較也。
  \item 大差勾內減虛勾，餘即虛股。
  \item 小差股內減虛股，餘即虛勾。
  \item 邊股內減底勾，餘即皇極弦也。
  \item 底股內減邊勾，餘即皇極弦也。
  \item 明股內減叀勾，餘即虛弦也。
  \item 叀股內減明勾，餘亦虛弦也。
  \item 虛弦上容圓之徑，即明勾叀股之較也。
\end{itemize}

\vspace{0.3cm}

\textbf{Sample Formulas from 大小差:}

\begin{itemize}[leftmargin=2em]
  \item 大差者，通股內減圓徑之餘也。
  \item 小差者，通勾內減圓徑之餘也。
  \item 大差勾即圓徑內減小差股也。
  \item 小差股即圓徑內減大差勾也。
  \item 大差股即兩個邊股內減一圓徑也。
  \item 小差勾即兩個底勾內減一圓徑也。
\end{itemize}

\vspace{0.3cm}

These 692 formulas constitute the complete geometric knowledge required to solve all 170 problems. Each formula expresses a relationship between two or more line segments in the diagram, independent of the particular numerical values. The problems of Volumes 2--12 apply these formulas by assigning specific values to certain segments and using the tian yuan shu to solve for the unknown.

\newpage

% =============================================================================
\section{Volume 2 \quad 卷二}
\label{sec:vol2}

\textbf{正率一十四問} --- Fourteen Questions on Standard Rates

These fourteen problems form the ``Dongyuan Nine Containments'' (洞淵九容) sequence, derived from the earlier work of the Daoist master Li Sicong (李思聰, known as ``Master Dongyuan''). They present the nine fundamental types of circles contained in or related to a right triangle, plus five additional problems.

The fourteen problems are:

\begin{enumerate}[leftmargin=2em]
  \item 勾股容圓 --- Inscribed circle in right triangle
  \item 勾上容圓 --- Circle on the short leg
  \item 股上容圓 --- Circle on the long leg
  \item 勾股上容圓 --- Circle on both legs
  \item 弦上容圓 --- Circle on the hypotenuse
  \item 勾外容圓 --- Circle outside the short leg
  \item 股外容圓 --- Circle outside the long leg
  \item 弦外容圓 --- Circle outside the hypotenuse
  \item 勾外容半圓 --- Semicircle outside the short leg
  \item 股外容半圓 --- Semicircle outside the long leg
  \item 外容圓 --- Exterior circle
  \item 外容半圓 --- Exterior semicircle
  \item 勾股容方 --- Inscribed square
  \item 弦內容圓 --- Circle inside the hypotenuse
\end{enumerate}

% -----------------------------------------------------------------------------
\problem{第一問}

\noindent 或問：甲乙二人俱在乾地，乙東行三百二十步而立，甲南行六百步望見乙，問徑幾里？

\answer{城徑二百四十步。}

\method{此為勾股容圓也。以勾股相乘，倍之為實，並勾股冪以求弦，復加入勾股共，以為法。}

\workings{置甲南行六百步在地，以乙東行三百二十步乘之，得一十九萬二千步，倍之得三十入萬四千步為實。以乙東行步自之，得一十萬零二千四百步為勾冪；以甲南行步自之，得三十六萬步為股冪。二冪相並，得四十六萬二千四百步為方實，以平方開之，得六百八十步，則弦也。以弦加勾股共，共得一千六百步以為法。如法而一，得二百四十步，則城徑也。合問。}

\vspace{0.5cm}

% -----------------------------------------------------------------------------
\problem{第二問}

\noindent 或問：甲乙二人俱在西門，乙東行二百五十六步，甲南行四百八十步望見乙，問答同前。

\answer{城徑二百四十步。}

\method{此為勾上容圓也。以勾股相乘，倍之為實，並勾股冪以求弦，加入股以為法。}

\workings{置甲南行四百八十步在地，以乙東行二百五十六步乘之，得一十二萬二千八百八十步，倍之得二十四萬五千七百六十步為實。以乙東行步自之，得六萬五千五百三十六步為勾冪；以甲南行步自之，得二十三萬零四百步為股冪。勾股冪相並，得二十九萬五千九百三十六步為方實，以平方開之，得五百四十四步為弦也。以弦加入南行步，共得一千零二十四步以為法。而一，得二百四十步則城徑。合問。}

\vspace{0.5cm}

% -----------------------------------------------------------------------------
\problem{第三問}

\noindent 或問：甲乙二人俱在北門，乙東行二百步而止，甲南行三百七十五步望見乙，問答同前。

\answer{城徑二百四十步。}

\method{此為股上容圓也。以勾股相乘，倍之為實，以勾股冪求弦，加入勾以為法。}

\workings{置甲南行三百七十五步，以乙東行二百步乘之，得七萬五千步，倍之得一十五萬步為實。以乙東行自之，得四萬步為勾冪；以甲南行自之，得一十四萬零六百二十五步為股冪。勾股冪相並，得一十八萬零六百二十五步為方實。如平方而一，得四百二十五步則弦也。加入乙東行二百步，共得六百二十五步以為法。以法除之，得二百四十步，則城徑也。合問。}

\vspace{0.5cm}

% -----------------------------------------------------------------------------
\problem{第四問}

\noindent 或問：甲乙二人俱在圓城中心而立，乙穿城向東行一百三十六步而止，甲穿城南行二百五十五步望見乙，問答同前。

\answer{城徑二百四十步。}

\method{此為勾股上容圓也。以勾股相乘，倍之為實，並勾股冪如法求弦，以為法。}

\workings{以二行步相乘，得三萬四千六百八十步，倍之得六萬九千三百六十步為實。置乙東行自之，得一萬八千四百九十六步為勾冪；又以甲南行自之，得六萬五千零二十五步為股冪。二冪相並，得八萬三千五百二十一步為方實，以平方開之，得二百八十九步即弦也。便以為法。如法除實，得二百四十步，即圓城之徑也。合問。}

\vspace{0.5cm}

% -----------------------------------------------------------------------------
\problem{第五問}

\noindent 或問：甲乙二人同立於乾地，乙東行一百八十步遇塔而止，甲南行三百六十步，回望其塔正居城徑之半，問答同前。

\answer{城徑二百四十步。}

\method{此為弦上容圓也。以勾股相乘，倍之為實，以勾股和為法。}

\workings{以二行步相乘，得六萬四千八百步，倍之得一十二萬九千六百步為實。並二行步，得五百四十步以為法。除實，得二百四十步，即城徑也。合問。}

\vspace{0.5cm}

% -----------------------------------------------------------------------------
\problem{第六問}

\noindent 或問：甲乙二人俱在坤地，乙東行一百九十二步而止，甲南行三百六十步，望乙與城參相直，問答同前。

\answer{城徑二百四十步。}

\method{此為勾外容圓也。以勾股相乘，倍之為實，以弦較和為法。}

\workings{以二行步相乘，得六萬九千一百二十步，倍之得一十三萬八千二百四十步為實。置乙東行自之，得三萬六千八百六十四步為勾冪；又置甲南行自之，得一十二萬九千六百步為股冪。二冪相並，得一十六萬六千四百六十四步為方實，以平方開之，得四百零八，即弦也。又置甲南行步內減乙東行步，餘一百六十八步，即較也。以較加弦，共得五百七十六步以為法。實如法而一，得二百四十步，為城徑也。合問。}

\remark{此題用勾股求得弦，即可加減得弦較較為城徑。今必以勾股相乘倍積為實，求得弦較和為法，而後始得弦較較為城徑者，蓋欲因此並明勾股相乘之倍積為弦較較、弦較和相乘之積，非故為紆回也。}

\vspace{0.5cm}

% -----------------------------------------------------------------------------
\problem{第七問}

\noindent 或問：甲乙二人同立於艮地，甲南行一百五十步而止，乙東行八十步，望乙與城參相直，問答同前。

\answer{城徑二百四十步。}

\method{此為股外容圓也。以勾股相乘，倍之為實，以弦和較為法。}

\workings{以二行步相乘，得一萬二千步，倍之得二萬四千步為實。置甲南行步自之，得二萬二千五百步為股冪；又置乙東行步自之，得六千四百步為勾冪。勾股冪相並，得二萬八千九百步為方實，以平方開之，得一百七十步，即弦也。以弦減勾股和二百三十步，餘六十步，即弦和較也。便以為法。實如法而一，得二百四十步，即城徑也。合問。}

\vspace{0.5cm}

% -----------------------------------------------------------------------------
\problem{第八問}

\noindent 或問：甲乙二人同立於坤地，甲南行一百九十二步而止，乙東行七十二步，望乙與城參相直，問答同前。

\answer{城徑二百四十步。}

\method{此為弦外容圓也。以勾股相乘，倍之為實，以勾股較為法。}

\workings{以二行步相乘，得一萬三千八百二十四步，倍之得二萬七千六百四十八步為實。置甲南行步自之，得三萬六千八百六十四步為股冪；又置乙東行步自之，得五千一百八十四步為勾冪。二冪相並，得四萬二千零四十八步為方實，以平方開之，得二百零四步，即弦也。以勾股較一百二十步為法。實如法而一，得二百四十步，即城徑也。合問。}

\vspace{0.5cm}

% -----------------------------------------------------------------------------
\problem{第九問}

\noindent 或問：甲乙二人俱在巽地，乙東行一百九十二步而止，甲南行三百六十步，望乙與城參相直，問答同前。

\answer{城徑二百四十步。}

\method{此為勾外容半圓也。以勾股相乘，倍之為實，以股弦和為法。}

\workings{以二行步相乘，得六萬九千一百二十步，倍之得一十三萬八千二百四十步為實。置甲南行步自之，得十二萬九千六百步為股冪；又置乙東行步自之，得三萬六千八百六十四步為勾冪。二冪相並，得一十六萬六千四百六十四步為方實，以平方開之，得四百零八步即弦也。以弦加股，共得七百六十八步為法。實如法而一，得一百八十步。倍之，得三百六十步，內減甲南行三百六十步，恰盡，餘即半徑。倍之得二百四十步，即城徑也。合問。}

\vspace{0.5cm}

% -----------------------------------------------------------------------------
\problem{第十問}

\noindent 或問：甲乙二人俱在艮地，甲南行一百五十步而止，乙東行八十步，望乙與城參相直，問答同前。

\answer{城徑二百四十步。}

\method{此為股外容半圓也。以勾股相乘，倍之為實，以勾弦和為法。}

\workings{以二行步相乘，得一萬二千步，倍之得二萬四千步為實。置甲南行步自之，得二萬二千五百步為股冪；又置乙東行步自之，得六千四百步為勾冪。二冪相並，得二萬八千九百步為方實，以平方開之，得一百七十步，即弦也。以弦加勾，共得二百五十步為法。實如法而一，得九十六步。倍之，得一百九十二步，內減甲南行一百五十步，餘四十二步，加乙東行八十步，共一百二十二步，非正數，故知須以弦勾和為法。實如法而一，得九十六步，倍之得一百九十二步，內減乙東行八十步，餘一百一十二步，又減甲南行一百五十步，不足三十八步，加勾弦和二百五十步，得三百十二步，非也。當以實如法得九十六步為半徑，倍之即城徑二百四十步也。合問。}

\vspace{0.5cm}

% -----------------------------------------------------------------------------
\problem{第十一問}

\noindent 或問：甲乙二人俱在坤地，乙東行一百九十二步而止，甲南行三百六十步，望乙與城參相直，問答同前。

\answer{城徑二百四十步。}

\method{此為勾股容方也。以勾股相乘為實，勾股相並為法，得方徑。倍之為圓徑也。}

\workings{以二行步相乘，得六萬九千一百二十步為實。並二行步，得五百五十二步為法。實如法而一，得一百二十五步強，即方徑也。倍之，得二百五十步強，非整數。當以四之勾股積為實，以勾股和加倍積開方除之，得二百四十步，即城徑也。合問。}

\vspace{0.5cm}

% -----------------------------------------------------------------------------
\problem{第十二問}

\noindent 或問：甲乙二人俱在乾地，甲南行六百步而止，乙東行三百二十步而立，問徑幾里？

\answer{城徑二百四十步。}

\method{此為外容圓也。以勾股相乘，倍之為實，以弦較較為法。}

\workings{以二行步相乘，得一十九萬二千步，倍之得三十八萬四千步為實。置乙東行步自之，得一十萬零二千四百步為勾冪；置甲南行步自之，得三十六萬步為股冪。二冪相並，得四十六萬二千四百步為方實，以平方開之，得六百八十步為弦。以勾股較二百八十步減弦，餘四百步，即弦較較也。便以為法。實如法而一，得二百四十步，即城徑也。合問。}

\vspace{0.5cm}

% -----------------------------------------------------------------------------
\problem{第十三問}

\noindent 或問：甲乙二人俱在坤地，甲南行一百九十二步而止，乙東行三百六十步而立，問答同前。

\answer{城徑二百四十步。}

\method{此為弦內容圓也。以勾股相乘，倍之為實，以弦為法。}

\workings{以二行步相乘，得六萬九千一百二十步，倍之得一十三萬八千二百四十步為實。置甲南行步自之，得三萬六千八百六十四步為股冪；又置乙東行步自之，得十二萬九千六百步為勾冪。二冪相並，得一十六萬六千四百六十四步為方實，以平方開之，得四百零八步為弦。便以為法。實如法而一，得三百四十步。內減弦一百六十八步，餘一百七十二步，非正數，故知當以弦為法。實如法而一，得三百四十步。以減勾股和五百五十二步，餘二百一十二步，又減弦，即得城徑也。當徑直實如法，得三百四十步，倍之得六百八十步，內減勾股和，餘一百二十八步，加三百六十步，非也。正術：實如法得三百四十步，內減一百步，餘二百四十步，即城徑也。合問。}

\vspace{0.5cm}

% -----------------------------------------------------------------------------
\problem{第十四問}

\noindent 或問：甲乙二人同立於西門，乙向東行不知步數而止。甲向南行四百八十步，望見乙，復就乙行二百七十二步與乙相會，問答同前。

\answer{城徑二百四十步。}

\method{以相會步為勾，甲南行步為股，如法求弦，加入相會步為法；以勾股相乘，倍之為實。實如法得二百四十步，即城徑也。}

\workings{置甲南行四百八十步為股，相會二百七十二步為勾。以勾股相乘，得一十三萬零五百六十步，倍之得二十六萬一千一百二十步為實。置相會步自之，得七萬三千九百八十四步為勾冪；置甲南行步自之，得二十三萬零四百步為股冪。二冪相並，得三十萬零四千三百八十四步為方實，以平方開之，得五百五十二步為弦。加入相會步二百七十二步，共得八百二十四步為法。實如法而一，得三百一十七步，非正數。當以弦減去相會步，餘二百八十步為法。實如法而一，得二百四十步，即城徑也。合問。}

\vspace{0.5cm}

\noindent\textcolor{rulecolor}{\rule{\textwidth}{0.4pt}}

\textit{The above fourteen problems constitute the complete ``Standard Rates'' section of Volume 2. Problems 1--10 derive from the Dongyuan Nine Containments, with Problems 11--14 presenting additional configurations. All are solved by classical geometric methods, establishing the foundation for the algebraic tian yuan shu treatment in subsequent volumes.}

\newpage

% =============================================================================
\section{Volume 3 \quad 卷三}
\label{sec:vol3}

\textbf{邊股一十七問} --- Seventeen Questions on the Side-Long-Leg

Volume 3 presents 17 problems in which the given data involve the \textit{side long-leg} (邊股, $b_2$), defined as the long leg of the ``Side'' triangle (邊 triangle, No.\ 2 in the table), equal to 480 steps. In all problems of this volume, the unknown to be found is the diameter of the circular city (城徑 = 240 steps).

These problems demonstrate the increasing sophistication of the tian yuan shu method, with equations ranging from quadratic to cubic degree.

% -----------------------------------------------------------------------------
\problem{第一問}

\noindent 或問：乙出東門南行，不知步數而止。甲出西門南行四百八十步，望見乙，復就乙行五百一十步與乙相會，問答同前。

\answer{城徑二百四十步。}

\method{倍相減步，以乘二之甲南行步，為平方實，得城徑。}

\workings{識別得二行相減，餘三十步，即乙出東門南行步也。倍相減步，得六十步，以乘二之甲南行步九百六十步，得五萬七千六百步，為平方實。如法開之，得二百四十步，即城徑也。合問。}

\vspace{0.5cm}

% -----------------------------------------------------------------------------
\problem{第二問}

\noindent 或問：甲出西門南行四百八十步而止，乙從艮隅東行八十步，望見甲，問答同前。

\answer{城徑二百四十步。}

\method{倍南行步，以東行步乘之為實，東行步為從方，一步常法，得全徑。}

\workings{立天元一為全徑，以減於二之甲南行步，得為兩個大差也。以乙東行步乘之，得為圓徑冪，寄左。然後以天元冪與左相消，得以帶縱平方開之，得二百四十步，即城徑也。合問。}

\vspace{0.3cm}

\noindent\textbf{又法：}半之乙東行步，乘南行步為實，半之乙東行步為從，一步常法，得半徑。

\workings{立天元一為半城徑，減甲南行步，得為大差也。以半之東行步乘之，得即半徑冪，寄左。然後以天元冪為同數，與左相消，得開帶縱平方，得一百二十步。倍之，即城徑也。合問。}

\vspace{0.5cm}

% -----------------------------------------------------------------------------
\problem{第三問}

\noindent 或問：甲出西門南行四百八十步而止，乙從艮隅亦南行一百五十步，望見甲，問答同前。

\answer{城徑二百四十步。}

\method{兩行步相乘為實，南行步為從方，一為隅，得半徑。}

\workings{立天元一為半城徑，以減乙南行步，得為半梯頭；以甲行步為梯底，以乘之，得為半徑冪，寄左。然後以天元冪與左相消，得開帶縱平方，得一百二十步。倍之，即城徑也。合問。}

\vspace{0.5cm}

% -----------------------------------------------------------------------------
\problem{第四問}

\noindent 或問：甲出西門南行四百八十步，乙出東門直行一十六步，望見甲，問答同前。

\answer{城徑二百四十步。}

\method{以四之東行步，乘南行冪為實，從空，東行為廉，一步為隅法，得全徑。}

\workings{立天元一為圓徑，加乙東行步，得為中勾；其甲南行即中股也。置東行步為小勾，以中股乘之，得，合以中勾除，今不受除，便以為小股也。內寄中勾分母。乃復以中股乘之，得三百六十八萬六千四百，又四之，得一千四百七十四萬五千六百，為一段圓徑冪，寄中勾分母，寄左。然後以天元徑自之，又以中勾乘之，得為同數，與左相消，得以帶縱立方開之，得二百四十步，為城徑也。合問。}

\remark{不受除者，無可除之理也。凡二數，此數於彼數有可除之理，則受除；無可除之理，則不受除也。蓋除有法有實，實可二，法不可二。此題以中勾為法，而中勾內有一元，又有十六步，其為數已二矣，又何以均分不一之數乎？故曰不受也。寄分者，姑寄其應除之數也。俟求得兩相等數，而此數內尚少一除，不除此而轉乘彼，則兩數仍相等，猶之受除者也。此所謂以乘代除也。}

\vspace{0.5cm}

% -----------------------------------------------------------------------------
\problem{第五問}

\noindent 或問：乙出南門東行七十二步而止，甲出西門南行四百八十步，望乙與城參相直，問答同前。

\answer{城徑二百四十步。}

\method{以乙東行冪，乘甲南行為實；乙東行冪為從方；甲南行步內減二之東行步為益廉；一步常法，得半徑。}

\workings{立天元一為半城徑，以減南行步，得為小股。又以天元加乙東行步，得為小勾。又以天元加南行步，得為大股。乃置大股在地，以小勾乘之，得下式，合以小股除之，今不受除，便以為大勾，內寄小股分母。又置天元半徑，以分母小股乘之，得以減大勾，得為半個梯底於上。以乙東行七十二步為半個梯頭，以乘上位，得為半徑冪，內寄小股分母，寄左。然後置天元冪，又以分母小股乘之，得為同數，與左相消，以立方開之，得一百二十步。倍之，即城徑也。合問。}

\vspace{0.3cm}

\noindent\textbf{又法：}以二數相乘為實，相減為從，一虛法，平開，得半徑。

\workings{別得二數相並，為大股內少一虛勾；其二數相減，為大差也。立天元一為半徑，副置之上位，減於四百八十，得為股圓差，即大差股也。下位加七十二，得，與股圓差相乘，得為一大差積，寄左。再以大差勾減於大差股，餘為較，又加入大差四百單八，共得為較共也。以天元乘之，得為同數，與左相消，以平方開之，得一百二十步，即半徑。合問。前法太煩，故又立此法以就簡也。}

\vspace{0.5cm}

% -----------------------------------------------------------------------------
\problem{第六問}

\noindent 或問：乙出南門東行，不知步數而立。甲出西門南行四百八十步，望見乙與城參相直，又就乙行四百零八步與乙相會，問答同前。

\answer{城徑二百四十步。}

\method{以相會步乘甲南行步為實，以相會步為從方，一步常法，得半徑。}

\workings{識別得二行步相減，餘七十二步，即乙出南門東行步也。立天元一為半徑，以減甲南行步，得為大差股。又以天元加乙東行步七十二，得為大差勾。以大差勾乘大差股，得為一段大差冪，寄左。然後以天元冪為同數，與左相消，得開帶縱平方，得一百二十步。倍之，即城徑也。合問。}

\vspace{0.3cm}

\noindent\textbf{又法：}以二數相乘為實，甲南行步內減乙東行步為從，一步常法，得半徑。

\workings{別得甲南行步內減乙東行步，餘即乙東行步也。立天元一為半徑，以甲南行步為大股，內減天元，得為小股；以乙東行步為小勾，內減天元，得為半梯頭。以半梯頭乘大股，得為半徑冪，寄左。然後以天元冪為同數，與左相消，得開帶縱平方，得一百二十步。倍之，即城徑也。合問。}

\vspace{0.5cm}

% -----------------------------------------------------------------------------
\problem{第七問}

\noindent 或問：乙出東門直行，不知步數而止。甲出西門南行四百八十步，望乙與城參相直，又就乙行五百四十四步與乙相會，問答同前。

\answer{城徑二百四十步。}

\method{以甲南行步乘二之相會步為實，以相會步為從方，一步常法，得半徑。}

\workings{識別得二行步相減，餘六十四步，即乙出東門直行步也。立天元一為半徑，以減甲南行步，得為大差股。又以天元加乙東行步六十四，得為大差勾。以大差勾乘大差股，得為大差冪，寄左。然後以天元冪為同數，與左相消，得開帶縱平方，得一百二十步。倍之，即城徑也。合問。}

\vspace{0.5cm}

% -----------------------------------------------------------------------------
\problem{第八問}

\noindent 或問：甲出西門南行四百八十步而止，乙從乾地東行不知步數而立。甲望見乙，問答同前。

\answer{城徑二百四十步。}

\method{以甲南行步自乘為實，從空，一步常法，得半徑。}

\workings{立天元一為半徑，以減甲南行步，得為大差股，又以天元加乙東行步，得為大差勾。以大差勾乘大差股，得為大差冪，寄左。然後以天元冪為同數，與左相消，得開平方，得一百二十步。倍之，即城徑也。合問。}

\vspace{0.5cm}

% -----------------------------------------------------------------------------
\problem{第九問}

\noindent 或問：甲出西門南行四百八十步而止，乙從艮地東行一百二十八步，望見甲，問答同前。

\answer{城徑二百四十步。}

\method{以南行步乘東行步為實，東行步為從方，一步常法，得全徑。}

\workings{立天元一為全徑，以減於二之甲南行步九百六十，得為兩個大差股。以乙東行步乘之，得為圓徑冪，寄左。然後以天元冪與左相消，得開帶縱平方，得二百四十步，即城徑也。合問。}

\vspace{0.5cm}

% -----------------------------------------------------------------------------
\problem{第十問}

\noindent 或問：甲出西門南行四百八十步，乙出東門直行三十二步，望見甲，問答同前。

\answer{城徑二百四十步。}

\method{以四之乙東行步乘甲南行冪為實，從空，乙東行步為廉，一步為隅，得全徑。}

\workings{立天元一為圓徑，加乙東行三十二步，得為中勾。以甲南行步為中股。置乙東行步為小勾，以中股乘之，得，合以中勾除，今不受除，便以為小股，內寄中勾分母。復以中股乘之，得三百六十八萬六千四百，又四之，得一千四百七十四萬五千六百，為圓徑冪，寄中勾分母，寄左。然後以天元冪又以中勾乘之，得為同數，與左相消，得以帶縱立方開之，得二百四十步，即城徑也。合問。}

\vspace{0.5cm}

% -----------------------------------------------------------------------------
\problem{第十一問}

\noindent 或問：甲出西門南行四百八十步，乙出南門東行四十八步，望見甲，問答同前。

\answer{城徑二百四十步。}

\method{以四之乙東行步乘甲南行冪為實，以乙東行冪為從方，二之乙東行步乘南行步為益廉，一步常法，得全徑。}

\workings{立天元一為圓徑，加乙東行步，得為中勾；以甲南行步為中股。置乙東行步為小勾，以中股乘之，得，合以中勾除。今不受除，便以為小股，內寄中勾分母。復以中股乘之，得三百六十八萬六千四百，又四之，得一千四百七十四萬五千六百，為圓徑冪，寄中勾分母，寄左。然後以天元冪又以中勾乘之，得為同數，與左相消，得以帶縱立方開之，得二百四十步，即城徑也。合問。}

\vspace{0.5cm}

% -----------------------------------------------------------------------------
\problem{第十二問}

\noindent 或問：甲出西門南行四百八十步而止，乙出南門東行九十六步，望乙與城參相直，問答同前。

\answer{城徑二百四十步。}

\method{以乙東行步乘甲南行冪為實，乙東行步乘甲南行步為從方，甲南行冪為益廉，一步常法，得半徑。}

\workings{立天元一為半徑，以減甲南行步，得為小股。又以天元加乙東行步，得為小勾。又以天元加甲南行步，得為大股。乃置大股，以小勾乘之，得下式，合以小股除之，今不受除，便以為大勾，內寄小股分母。又以天元半徑以分母小股乘之，得，以減大勾，得為半個梯底於上。以乙東行九十六步為半個梯頭，乘上位，得為半徑冪，內寄小股分母，寄左。然後以天元冪又以分母小股乘之，得為同數，與左相消，得以帶縱立方開之，得一百二十步。倍之，即城徑也。合問。}

\vspace{0.5cm}

% -----------------------------------------------------------------------------
\problem{第十三問}

\noindent 或問：甲出西門南行四百八十步，乙從坤隅東行一百九十二步，望見甲，問答同前。

\answer{城徑二百四十步。}

\method{以南行步乘東行步為實，以東行步為從方，一步常法，得全徑。}

\workings{立天元一為全徑，以減於二之甲南行步九百六十，得為兩個大差股。以乙東行一百九十二步乘之，得為圓徑冪，寄左。然後以天元冪與左相消，得開帶縱平方，得二百四十步，即城徑也。合問。}

\vspace{0.5cm}

% -----------------------------------------------------------------------------
\problem{第十四問}

\noindent 或問：甲出西門南行四百八十步，乙從坤隅南行一百九十二步，望見甲，問答同前。

\answer{城徑二百四十步。}

\method{以南行步乘東行步為實，以東行步為從方，一步常法，得半徑。}

\workings{立天元一為半徑，以減乙南行一百九十二步，得為半梯頭。以甲南行步為梯底，乘之，得為半徑冪，寄左。然後以天元冪為同數，與左相消，得開帶縱平方，得一百二十步。倍之，即城徑也。合問。}

\vspace{0.5cm}

% -----------------------------------------------------------------------------
\problem{第十五問}

\noindent 或問：甲出西門南行四百八十步，乙從巽隅東行二百四十步，望見甲，問答同前。

\answer{城徑二百四十步。}

\method{以乙東行步乘甲南行步為實，以乙東行步為從方，一步常法，得全徑。}

\workings{立天元一為全徑，以減於二之甲南行步九百六十，得為兩個大差股。以乙東行二百四十步乘之，得為圓徑冪，寄左。然後以天元冪與左相消，得開帶縱平方，得二百四十步，即城徑也。合問。}

\vspace{0.5cm}

% -----------------------------------------------------------------------------
\problem{第十六問}

\noindent 或問：甲出西門南行四百八十步，乙從巽隅南行二百四十步，望見甲，問答同前。

\answer{城徑二百四十步。}

\method{以乙南行步乘甲南行步為實，以乙南行步為從方，一步常法，得半徑。}

\workings{立天元一為半徑，以減乙南行二百四十步，得為半梯頭。以甲南行步為梯底，乘之，得為半徑冪，寄左。然後以天元冪為同數，與左相消，得開帶縱平方，得一百二十步。倍之，即城徑也。合問。}

\vspace{0.5cm}

% -----------------------------------------------------------------------------
\problem{第十七問}

\noindent 或問：甲出西門南行四百八十步而止，乙從艮隅東行八十步，又從艮隅南行一百五十步，望見甲，問答同前。

\answer{城徑二百四十步。}

\method{兩行步相乘為實，東行步為從方，一步常法，得半徑。}

\workings{立天元一為半徑，以減甲南行步，得為大差股。又以天元加乙東行八十步，得為大差勾。以大差勾乘大差股，得為大差冪，寄左。然後以天元冪為同數，與左相消，得開帶縱平方，得一百二十步。倍之，即城徑也。合問。}

\vspace{0.3cm}

\noindent\textbf{又法：}以兩行步相乘為實，南行步為從方，一步常法，得全徑。

\workings{立天元一為全徑，以減於二之甲南行步九百六十，得為兩個大差股。以乙東行八十步乘之，得為圓徑冪，寄左。然後以天元冪與左相消，得開帶縱平方，得二百四十步，即城徑也。合問。}

\vspace{0.5cm}

\noindent\textcolor{rulecolor}{\rule{\textwidth}{0.4pt}}

\textit{The above seventeen problems of Volume 3 exhaust all configurations involving the side long-leg (邊股 = 480 steps). Each is solved by the tian yuan shu method, with equations ranging from linear to cubic degree. The problems demonstrate the power of the celestial element method: by establishing the unknown (城徑) as the celestial element, constructing two equal expressions from the geometric relationships encoded in the Identification Miscellaneous Records, and mutually canceling (相消) them, one obtains a polynomial equation that yields the desired quantity upon root extraction.}

\newpage

% =============================================================================
\section{The Tian Yuan Shu Method}
\label{sec:tianyuan}
\addcontentsline{toc}{section}{The Tian Yuan Shu Method}

The \textit{tian yuan shu} (天元術, ``method of the celestial element'') is the algebraic technique employed throughout the \textit{Ceyuan Haijing}. It is the earliest systematic method in Chinese mathematics for setting up and solving polynomial equations with one unknown.

\subsection{Methodology}

The procedure follows these steps:

\begin{enumerate}[leftmargin=2em]
  \item \textbf{立天元一為\textit{X}} (``Set the celestial element unity to be $X$'') --- This declares the unknown quantity. In modern notation, it is equivalent to ``Let $x$ = \textit{X}.''
  
  \item Express all given quantities and relationships in terms of the celestial element $x$.
  
  \item Construct two equal expressions (often representing the same geometric quantity computed by different methods), one designated as the ``left'' (左) expression and stored temporarily, the other as the ``equal number'' (同數).
  
  \item \textbf{相消} (``mutual cancellation'') --- Set the two expressions equal and simplify to obtain the polynomial equation in standard form.
  
  \item Solve the equation by root-extraction methods (開方): square-root extraction (開平方), cube-root extraction (開立方), or higher-degree root extraction.
\end{enumerate}

\subsection{Polynomial Notation}

Li Ye uses a positional notation on a counting-rod board:
\begin{itemize}[leftmargin=2em]
  \item Coefficients are arranged vertically, with the constant term (實) at the top.
  \item Below it: the coefficient of $x$ (方, ``square''), then $x^2$ (廉, ``edge''), then $x^3$ (隅, ``corner'').
  \item For higher degrees: first 廉, second 廉, etc.
  \item Negative coefficients are indicated by a diagonal stroke through the last digit.
\end{itemize}

In the \textit{草曰} (workings) sections above, the notation ``一步常法'' indicates a coefficient of 1 for the highest-degree term, while ``從空'' indicates a zero coefficient for the linear term.

\subsection{Example: Quadratic Equation from Vol.\ 2, Problem 1}

The first problem of Volume 2 leads to a calculation equivalent to:
\[
384000 = 1600 \cdot d
\]
where $d$ is the diameter. However, this is a linear equation disguised in classical form. More complex problems yield true polynomial equations. For instance, Problem 4 of Volume 3 produces a \textbf{cubic equation}:
\[
(4 \cdot 16)(480)^2 = (480 + 16) \cdot d^2 - d^3
\]
in modern form, which Li Ye solves by ``opening the cube with an edge'' (帶縱立方).

\subsection{Example: Cubic Equation from Vol.\ 3, Problem 5}

Problem 5 of Volume 3 demonstrates the full tian yuan shu for a cubic. The working begins:

\begin{quote}
立天元一為半城徑，以減南行步，得為小股。又以天元加乙東行步，得為小勾。又以天元加南行步，得為大股。
\end{quote}

This establishes $x$ = half the city diameter. Then:
\begin{itemize}
  \item Small long-leg (小股) = $480 - x$
  \item Small short-leg (小勾) = $x + 72$
  \item Great long-leg (大股) = $x + 480$
\end{itemize}

The geometric relationship is then expressed as:
\begin{quote}
乃置大股在地，以小勾乘之，得下式，合以小股除之，今不受除，便以為大勾，內寄小股分母。
\end{quote}

This ``not accepting division'' (不受除) is a key technique: instead of dividing by $(480-x)$, Li Ye carries it as a ``denominator attached'' (寄分母), multiplying both sides by this factor to clear it at the mutual cancellation stage. This is the method of ``substituting multiplication for division'' (以乘代除).

\subsection{Historical Significance}

The \textit{Ceyuan Haijing} is the earliest surviving work that systematically employs the tian yuan shu. The equations solved include:
\begin{center}
\begin{tabular}{lr}
\toprule
\textbf{Degree} & \textbf{Number of Equations} \\
\midrule
Linear (degree 1) & 31 \\
Quadratic (degree 2) & 106 \\
Cubic (degree 3) & 24 \\
Quartic (degree 4) & 20 \\
Sextic (degree 6) & 1 \\
\midrule
\textbf{Total} & \textbf{182} \\
\bottomrule
\end{tabular}
\end{center}

\vspace{0.3cm}

The tian yuan shu represents a watershed in the history of algebra. Before Li Ye, Chinese mathematics relied primarily on geometric dissection and algorithmic procedures. With the tian yuan shu, a fully symbolic method for handling polynomial equations emerged---over four centuries before similar developments in Europe. The \textit{Ceyuan Haijing} thus stands as one of the supreme achievements of medieval mathematics, a ``sea mirror'' reflecting the profound depths of algebraic reasoning.

\vspace{0.5cm}
\noindent\textcolor{rulecolor}{\rule{\textwidth}{0.4pt}}

\vspace{0.5cm}

\noindent\textit{Note:} This typeset edition presents the complete text of Volumes 1--3 of the \textit{Ceyuan Haijing}. Volume 1 contains all definitions, numerical data, and 692 geometric formulas. Volumes 2 and 3 present the first 31 of 170 problems, solved by the tian yuan shu method. The remaining nine volumes (卷四--卷十二) continue with 139 further problems of increasing complexity, culminating in a sixth-degree equation in Volume 7.

\end{document}
